% Generated draft: requires mathematical and source review.
\documentclass[11pt]{article}
\usepackage[margin=25mm]{geometry}
\usepackage{fontspec}
\setmainfont{Latin Modern Roman}
\usepackage{amsmath,amssymb,mathtools,unicode-math}
\setmathfont{Latin Modern Math}
\usepackage{enumitem,needspace,xurl,hyperref}
\setlength{\parindent}{0pt}
\setlength{\parskip}{5pt}
\setlength{\emergencystretch}{3em}
\newcommand{\N}{\mathbb{N}}
\newcommand{\Z}{\mathbb{Z}}
\newcommand{\Q}{\mathbb{Q}}
\newcommand{\R}{\mathbb{R}}
\newcommand{\C}{\mathbb{C}}
\newcommand{\F}{\mathbb{F}}
\newcommand{\A}{\mathbb{A}}
\newcommand{\PP}{\mathbb{P}}
\newcommand{\E}{\mathbb{E}}
\newcommand{\eps}{\varepsilon}
\newcommand{\dd}{\,\mathrm{d}}
\begin{document}
\section*{UnsolvedMath ID: 30004416}
\textbf{Source problem code:} catalogue entry
\subsection{Definitions and mathematical statement}
Let \(G\) be a finite group, and let \(\operatorname{Irr}(G)\) denote the set of characters of irreducible finite-dimensional complex representations of \(G\). For \(\chi\in\operatorname{Irr}(G)\), its degree is \(\chi(1)\). Define
\[
\operatorname{Irr}_{2'}(G)=\{\chi\in\operatorname{Irr}(G):2\nmid\chi(1)\},
\qquad
\Q(\chi)=\Q\bigl(\chi(g):g\in G\bigr)\subseteq\C.
\]
Thus \(\Q(\chi)\) is the field generated over \(\Q\) by all values of \(\chi\).

Determine the collection of field extensions
\[
F_2=\left\{\Q(\chi)/\Q:\ G\text{ is a finite group},\ \chi\in\operatorname{Irr}_{2'}(G)\right\}.
\]
Here \(\Q(\chi)/\Q\) denotes a field extension, not a quotient. The quantifiers range over all finite groups and all their irreducible complex characters of odd degree.

The supplied dated curation states that the classification for \(p=2\) has been obtained, although the registry labels this record partially solved. Its further discussion of conjectural classifications for odd primes concerns a broader problem and does not change the target \(F_2\). No classification theorem is supplied in the record, so none is asserted here.
\subsection{Short English statement}
Determine exactly which fields generated by character values occur for irreducible complex characters of odd degree, allowing every finite group. The supplied curation reports that this classification is known for \(p=2\); that literature assessment has not been independently verified.
\subsection{Sources}
\begin{itemize}
\item Local UnsolvedMath registry, supplied complete record, permanent ID \(30004416\), source problem code \texttt{catalogue entry}, title \emph{Character Fields of Odd-Degree Representations}. The record gives proposed year 2020 but names no proposer.
\item Supplied workshop citation: \emph{Mini-Workshop: Kronecker, Plethysm, and Sylow Branching Coefficients and their Applications to Complexity Theory} (2021), Oberwolfach Reports, pp.~517--568, DOI \texttt{10.4171/owr/2020/9}. \url{https://doi.org/10.4171/owr/2020/9}. The supplied literature list also labels this reference \emph{Oberwolfach Report 9/2020}.
\item Supplied literature reference, attributed to Hung, Navarro, and Tiep and labeled \emph{fields of values}: \url{https://arxiv.org/abs/2601.15987}.
\item Supplied literature reference labeled \emph{Fields of values of p-prime-degree characters}: \url{https://www.cambridge.org/core/journals/forum-of-mathematics-pi/article/fields-of-values-of-characters-of-degree-not-divisible-by-p/0FE327DAE6088EC9B43A0CFAFC8E8D52}.
\item Literature assessment checked at \texttt{2026-08-22}, according to the supplied record: fields of values of characters of degree prime to \(p\) are classified for \(p=2\), while the analogous classification for odd primes remains conjectural with partial results. This is supplied dated curation derived from the OWR material, not an independently verified assessment; its time-sensitive status and cited claims require independent verification.
\end{itemize}
\end{document}
